\section{Conclusion}
\label{sec:conclusion}

We tested whether LLMs take loopholes when they do not know an answer but still feel pressed to respond. To do so we constructed \kcl, the first paired-prompt benchmark for knowledge-conflict loophole behavior, and evaluated five LLMs spanning frontier, large-open, small-open, and reasoning tiers. Across all five models, the within-pair unsatisfactory rate is higher on unknowable than knowable items; a binomial GLM with cluster-robust SEs gives an odds ratio of $3.24$ ($p = 0.020$) for knowability, and the long-chain-of-thought \rseventy roughly halves the unsatisfactory rate relative to the small-open baseline (OR $0.50$, $p = 0.012$). The dominant mechanism, however, is confident hallucination, not the Choi-style verbalized reinterpretation we initially predicted. A complementary \selfaware probe inverts the picture: on genuinely unanswerable items, models almost universally produce verbal dodges and almost never refuse.

The headline takeaway is that the original ``LLMs take loopholes when they do not know'' hypothesis is correct but mechanistically two-modal. On factual questions the loophole is implicit---it lives in the model's choice to assert rather than abstain. On underspecified questions the loophole is verbal---a non-committal essay that addresses related considerations. The two modes share the substrate (the model wants to say something) but the surface signature is different, which is why a single refusal-vs-answer score misses both.

\para{Future work.} Four follow-ups are immediate. (i) A second judge (\haiku) on the same 600 \kcl generations, with reported Cohen's $\kappa$, would close the largest internal-validity threat. (ii) Internal-state probes \citep{azaria2023internal, orgad2024llmsknow, kossen2024semantic} on open-weights models during a confident-wrong \kcl response would directly test the implicit-loophole account: if the internal probe fires for ``I do not know'' while the surface asserts a number, the decision-stage interpretation is confirmed. (iii) A refusal-encouraging system-prompt variant analogous to \citet{choi2025loopholes}'s \texttt{scalar\_max} would measure how closeable the gap is with a one-line nudge. (iv) Scaling \kcl to $\ge 200$ templates would give per-category power and let us ask whether the implicit-loophole rate differs across \texttt{precise-numeric}, \texttt{private-people}, and \texttt{count-or-list}. Multi-turn dynamics---whether a follow-up ``are you sure?'' flips the model from confident-wrong to refusal---is a fifth, broader direction.

The benchmark, judging rubric, generation cache, and analysis scripts are open and reproducible from a single \texttt{pyproject.toml}.
